% ===========================================================================
% SOURCES: docs/references/*.pdf (the papers are already in the repository)
%          docs/controller_design.md sections 1-2 (the Angel-Eye / WP486 study
%          that chose this architecture) - in Greek, translate as you go
%
% WRITE THIS AFTER the design chapters. "Positioning" is a claim about what you
% built, so it cannot be written before you have written what you built.
% ===========================================================================
\chapter{Related Work}
\label{chap:related}

\section{Deep Learning for Plant Disease Detection}
\label{sec:related-plant}
% PlantVillage [hughes2015plantvillage] and the CNN baseline
% [mohanty2016plantdisease]; the grapevine IoT service [morellos2024iot].
% The point to land: the accuracy problem is largely solved in software, and
% the open problem is doing it in the field within a power budget. That is the
% gap this thesis addresses, and it is why the thesis is an architecture
% contribution rather than a modelling one.
%
% Be explicit about the dataset-leakage issue here as well as in
% Section~\ref{sec:leakage}: published accuracies on PlantVillage-derived splits are
% frequently inflated by images of the same leaf on both sides of the split -
% augmented copies, and, less visibly, the several photographs PlantVillage
% took of every leaf. [mohanty2016plantdisease] kept each leaf's photos on one
% side and published the leaf grouping that makes this possible; a split by
% file or by image id ignores it. This thesis found both kinds in its own
% splits (DD-034) and now splits by leaf using that grouping.

\section{CNN Accelerators on FPGAs}
\label{sec:related-accel}
% Organise by what each one decided, not chronologically:
%   - loop tiling / design-space exploration: [zhang2015optimizing],
%     [ma2017optimizing]
%   - embedded end-to-end flows: [qiu2016goingdeeper], [guo2018angeleye]
%   - framework generators: [umuroglu2017finn], [venieris2016fpgaconvnet],
%     [duarte2018hls4ml]
%   - depthwise-separable specifically: [bai2018dwsep]
%   - dataflow taxonomy and the energy argument: [chen2017eyeriss],
%     [sze2017efficient]
%   - the vendor answer: the Xilinx DPU [xilinx_pg338]
% Angel-Eye [guo2018angeleye] deserves the most space: it is the closest
% comparable (embedded Zynq, quantised, instruction-driven) and it is what
% controller_design.md actually studied before choosing this architecture.

\section{Positioning of This Thesis}
\label{sec:positioning}
% Three honest distinctions, none of them "ours is faster":
%
%   1. VERIFICATION. Most of the accelerator literature reports accuracy after
%      quantisation. This design reports bit-exactness against its own software
%      reference - \thElements elements over \thInstructions instructions - which
%      is a strictly stronger and much rarer claim. Chapter~\ref{chap:verification}.
%
%   2. HAND-WRITTEN RTL over HLS [xilinx_pg338, venieris2016fpgaconvnet]. Say
%      plainly that this costs development time and buys control; cite the
%      timing work in Chapter~\ref{chap:timing} as the evidence that the control was
%      worth having, since the 4.6x Fmax improvement came from decisions at a
%      level HLS does not expose.
%
%   3. SCOPE. A single network, fully specialised, rather than a general
%      framework. The instruction set exists (Chapter~\ref{appendix:isa}) but it compiles one
%      model. Say so; it is a limitation, not a feature.
